\begin{helper}[Defective original-tree confinement]
Let $L$ be a \noderef{ScheideLastExitPath} certificate for $\alpha$.
Suppose that its last inside vertex $v$ belongs to
$X_0=\noderef{TreeSequenceVertexSet}(T)$ and that
$\alpha\in\noderef{DefectiveColors}(G,\varphi,X_0)$.  Regard $\alpha$
as the corresponding defective index $d$, put
$U_\alpha=\{(R d).u1,(R d).u2\}$, and put
$X_\alpha=X_0\cup U_\alpha$.  Then $D.f0(d)$ is the unique
$\alpha$-colored boundary identity of $X_\alpha$, every vertex-node of
$L$'s path from $y_0$ through $v$ lies in $X_\alpha$, and every edge
identity belonging both to $\noderef{BoundaryEdges}(G,X)$ and to the full
$(D.\alpha_0,\alpha)$ component through $y_0$ is the crossing edge of
$L$.
\end{helper}
\begin{proof}
Let $i=D.\gammaIndex(d)$ and
$U_\gamma=\{C.\mathit{blockVertex}(i,0),C.\mathit{blockVertex}(i,1)\}$.
The extension field of \noderef{ScheideReducedConfiguration} supplies the
inherited \noderef{ScheideStructuralConfiguration}.  Its exact block
positions, together with the image definition
\noderef{TreeSequenceVertexSet}, give $U_\gamma\subseteq X_0$.
The first conclusion of \noderef{ScheideSymmetricDefectiveInternal}
supplies an $\alpha$-colored edge internal to $U_\gamma$.  Its endpoints
are the two distinct members of $U_\gamma$ by
\noderef{LooplessMultigraph}.  The original coloring is proper because
that structural configuration contains a
\noderef{GlobalMaximumTriple}, hence a \noderef{KTriple} and an
\noderef{IsProperEdgeColoring}.  Properness and
\noderef{AdjacentEdges} therefore exclude every other $\alpha$-edge
incident with either deleted vertex: by \noderef{Incident}, any such edge
and the internal edge would be distinct edges incident with a common
vertex, hence adjacent.

Consequently every $\alpha$-colored boundary edge of $X_0$ has its
tree endpoint in $X_0\setminus U_\gamma$.  Indeed,
\noderef{BoundaryEdges} gives it exactly one endpoint in $X_0$, and the
preceding uniqueness of $\alpha$-incidence excludes that endpoint from
$U_\gamma$.  The exact-three field of
\noderef{ScheideReplacementData} now identifies all such identities as
$D.f0(d),(R d).f1,(R d).f2$.  The latter two have, respectively, one
endpoint in the retained set and the endpoints $(R d).u1,(R d).u2$ in
$U_\alpha$; by \noderef{Incident}, they therefore join those vertex
identities and are internal to $X_\alpha$.  They are distinct and their
$U_\alpha$ endpoints are distinct: otherwise the endpoint records and
\noderef{Incident} would make two distinct edges of the same color
incident with a common vertex and hence adjacent, contradicting
properness.  Thus every vertex of $U_\alpha$ is already incident with one
of these two $\alpha$-edges.  Properness excludes any further
$\alpha$-incidence there, so in particular
\noderef{BoundaryEdges} excludes any further $\alpha$-edge leaving
$X_\alpha$ through $U_\alpha$.

The outside endpoint of $D.f0(d)$ is outside $X_0$ by
\noderef{ScheideDefectiveSetup}.  It is not either member of
$U_\alpha$, since equality would make $D.f0(d)$ adjacent to the
corresponding distinct edge $(R d).f1$ or $(R d).f2$ of the same color:
the endpoint records and \noderef{Incident} give the common incidence.
Its other endpoint is the tree vertex recorded in
\noderef{ScheideDefectiveSetup}, hence lies in $X_0$, so
\noderef{BoundaryEdges} says that $D.f0(d)$ leaves $X_\alpha$.  The
preceding exhaustive list and boundary orientation prove that it is the
unique member of
\noderef{BoundaryColorEdges}$(G,\varphi,X_\alpha,\alpha)$.

The set $X_\alpha$ is closed by
\noderef{ScheideReplacementUnionClosed}.  The first base position in
\noderef{ScheideStructuralConfiguration} and the definition
\noderef{TreeSequenceVertexSet} put $y_0$ in $X_0\subseteq X_\alpha$, and
\noderef{ScheideDefectiveSetup} says that $y_0$ misses $\alpha_0$.  Hence
\noderef{ColoredVertexSetClosed} forbids an
$\alpha_0$-colored boundary edge of $X_\alpha$.  The crossing $g$ of
$L$ has one endpoint $v\in X_0\subseteq X_\alpha$ and its other endpoint
outside $X$.  The union formula in
\noderef{ScheideEnlargedVertexSet} gives $X_\alpha\subseteq X$, so this
second endpoint is outside $X_\alpha$.  Thus
\noderef{BoundaryEdges} makes $g$ a boundary edge of $X_\alpha$; it has
color $\alpha$, so the uniqueness just proved gives $g=D.f0(d)$.

If the path $P$ of $L$ left $X_\alpha$ before reaching $v$, choose its
first departure.  The adjacent vertex-nodes lie respectively inside and
outside $X_\alpha$, so \noderef{BoundaryEdges} makes the intervening
edge-node a boundary identity.  Its edge has color $\alpha_0$ or $\alpha$
by \noderef{TwoColorIncidenceGraph}; closedness excludes the first color,
and the second makes it the unique edge $D.f0(d)=g$.  But this earlier
edge-node is distinct from the edge-node $g$ after $v$, contradicting the
simplicity recorded in \noderef{ScheideLastExitPath}.  Thus all
vertex-nodes through $v$ lie in $X_\alpha$.

Finally let $f$ be any boundary edge of $X$ in the full component.  The
support equality for $L$ puts $f$ on $P$.  Before and through $v$, every
vertex-node is in $X_\alpha\subseteq X$; after $v$, every vertex-node is
outside $X$.  By \noderef{BoundaryEdges}, $f$ has exactly one endpoint in
$X$; by \noderef{Incident} and \noderef{TwoColorIncidenceGraph}, its two
endpoint vertex-nodes are its two neighbors on the component path.  Hence
the only possible path edge-node with one endpoint vertex in $X$ and the
other outside is the edge immediately after $v$, namely $g$.  Therefore
$f=g$.
\end{proof}
